% GENERATED FILE. Edit canonical lessons, then run npm run generate:books.
% canonical-source-hash: fnv1a64:e789aac7cf8e78c2
% canonical-lessons: HI-C243-nakam, HI-C243-ghamandi, HI-C243-sidhasada, HI-C243-calak, HI-C243-jhagralu

\chapter{Unsuccessful, Proud, Simple, Cunning, Quarrelsome}
\label{ch:hi-unsuccessful-proud-simple-cunning-quarrelsome}
\hlchaptermodality{243}

\begin{chapteropening}
\textbf{By the end of this chapter:} \emph{I can say unsuccessful, proud, simple, cunning and quarrelsome.}

Complete the last lesson of chapter 243: I can say unsuccessful, proud, simple, cunning and quarrelsome.
\end{chapteropening}

\section[\texorpdfstring{\hi{नाकाम} (nākām)}{नाकाम (nākām)}]{\hi{नाकाम} (nākām) — unsuccessful}

\label{lesson:HI-C243-nakam}



\begin{quote}
Before the new one: say the Hindi for useless, then the Hindi for successful.
\end{quote}

\subsection*{You'll want to know: \hi{नाकाम}}
\textbf{\hi{नाकाम}} — \emph{nākām} — \textquotedblleft{}unsuccessful\textquotedblright{}.

\begin{grammarlens}[title={What you've built}]
One more word for this chapter.
\end{grammarlens}

\subsection*{Your turn}
\begin{itemize}
\raggedright
  \item \emph{Say it:} \emph{nākām}
  \item \emph{Say it:} \emph{nākām}, once more
  \item \emph{Say it:} say \emph{nākām}
  \item \emph{Recall:} say the Hindi for useless, then the Hindi for successful, then say \emph{nākām} again
\end{itemize}

\begin{tcolorbox}[breakable,colback=teal!4,colframe=teal!35!black,title={Before you move on}]
What does \hi{नाकाम} mean? (\textquotedblleft{}Unsuccessful\textquotedblright{}.) Say it once more.
\end{tcolorbox}
\section[\texorpdfstring{\hi{घमंडी} (ghamaṇḍī)}{घमंडी (ghamaṇḍī)}]{\hi{घमंडी} (ghamaṇḍī) — proud, arrogant}

\label{lesson:HI-C243-ghamandi}



\begin{quote}
Before the new one: say the Hindi for successful, then the Hindi for unsuccessful.
\end{quote}

\subsection*{You'll want to know: \hi{घमंडी}}
\textbf{\hi{घमंडी}} — \emph{ghamaṇḍī} — \textquotedblleft{}proud, arrogant\textquotedblright{}.

\begin{grammarlens}[title={What you've built}]
One more word for this chapter.
\end{grammarlens}

\subsection*{Your turn}
\begin{itemize}
\raggedright
  \item \emph{Say it:} \emph{ghamaṇḍī}
  \item \emph{Say it:} \emph{ghamaṇḍī}, once more
  \item \emph{Say it:} say \emph{ghamaṇḍī}
  \item \emph{Recall:} say the Hindi for successful, then the Hindi for unsuccessful, then say \emph{ghamaṇḍī} again
\end{itemize}

\begin{tcolorbox}[breakable,colback=teal!4,colframe=teal!35!black,title={Before you move on}]
What does \hi{घमंडी} mean? (\textquotedblleft{}Proud, arrogant\textquotedblright{}.) Say it once more.
\end{tcolorbox}
\section[\texorpdfstring{\hi{सीधासादा} (sīdhāsādā)}{सीधासादा (sīdhāsādā)}]{\hi{सीधासादा} (sīdhāsādā) — simple, unpretentious}

\label{lesson:HI-C243-sidhasada}



\begin{quote}
Before the new one: say the Hindi for unsuccessful, then the Hindi for proud.
\end{quote}

\subsection*{You'll want to know: \hi{सीधासादा}}
\textbf{\hi{सीधासादा}} — \emph{sīdhāsādā} — \textquotedblleft{}simple, unpretentious\textquotedblright{}.

\begin{grammarlens}[title={What you've built}]
One more word for this chapter.
\end{grammarlens}

\subsection*{Your turn}
\begin{itemize}
\raggedright
  \item \emph{Say it:} \emph{sīdhāsādā}
  \item \emph{Say it:} \emph{sīdhāsādā}, once more
  \item \emph{Say it:} say \emph{sīdhāsādā}
  \item \emph{Recall:} say the Hindi for unsuccessful, then the Hindi for proud, then say \emph{sīdhāsādā} again
\end{itemize}

\begin{tcolorbox}[breakable,colback=teal!4,colframe=teal!35!black,title={Before you move on}]
What does \hi{सीधासादा} mean? (\textquotedblleft{}Simple, unpretentious\textquotedblright{}.) Say it once more.
\end{tcolorbox}
\section[\texorpdfstring{\hi{चालाक} (cālāk)}{चालाक (cālāk)}]{\hi{चालाक} (cālāk) — cunning, clever}

\label{lesson:HI-C243-calak}



\begin{quote}
Before the new one: say the Hindi for proud, then the Hindi for simple.
\end{quote}

\subsection*{You'll want to know: \hi{चालाक}}
\textbf{\hi{चालाक}} — \emph{cālāk} — \textquotedblleft{}cunning, clever\textquotedblright{}.

\begin{grammarlens}[title={What you've built}]
One more word for this chapter.
\end{grammarlens}

\subsection*{Your turn}
\begin{itemize}
\raggedright
  \item \emph{Say it:} \emph{cālāk}
  \item \emph{Say it:} \emph{cālāk}, once more
  \item \emph{Say it:} say \emph{cālāk}
  \item \emph{Recall:} say the Hindi for proud, then the Hindi for simple, then say \emph{cālāk} again
\end{itemize}

\begin{tcolorbox}[breakable,colback=teal!4,colframe=teal!35!black,title={Before you move on}]
What does \hi{चालाक} mean? (\textquotedblleft{}Cunning, clever\textquotedblright{}.) Say it once more.
\end{tcolorbox}
\section[\texorpdfstring{\hi{झगड़ालू} (jhagṛālū)}{झगड़ालू (jhagṛālū)}]{\hi{झगड़ालू} (jhagṛālū) — quarrelsome}

\label{lesson:HI-C243-jhagralu}



\begin{quote}
Before the new one: say the Hindi for simple, then the Hindi for cunning.
\end{quote}

\subsection*{You'll want to know: \hi{झगड़ालू}}
\textbf{\hi{झगड़ालू}} — \emph{jhagṛālū} — \textquotedblleft{}quarrelsome\textquotedblright{}.

\begin{grammarlens}[title={What you've built}]
One more word for this chapter.
\end{grammarlens}

\subsection*{Your turn}
\begin{itemize}
\raggedright
  \item \emph{Say it:} \emph{jhagṛālū}
  \item \emph{Say it:} \emph{jhagṛālū}, once more
  \item \emph{Say it:} say \emph{jhagṛālū}
  \item \emph{Recall:} say the Hindi for simple, then the Hindi for cunning, then say \emph{jhagṛālū} again
\end{itemize}

\begin{tcolorbox}[breakable,colback=teal!4,colframe=teal!35!black,title={Before you move on}]
What does \hi{झगड़ालू} mean? (\textquotedblleft{}Quarrelsome\textquotedblright{}.) Say it once more.
\end{tcolorbox}
